\documentclass[12pt]{article}
\usepackage[margin=1in]{geometry}
\usepackage{listings}
\usepackage{xcolor}
\usepackage{enumitem}
\usepackage{booktabs}
\usepackage{array}

\definecolor{codebg}{gray}{0.95}
\definecolor{codetext}{rgb}{0.1,0.1,0.5}

\lstset{
  backgroundcolor=\color{codebg},
  basicstyle=\ttfamily\small\color{codetext},
  frame=single,
  framerule=0pt,
  rulecolor=\color{codebg},
  breaklines=true,
  xleftmargin=8pt,
  xrightmargin=8pt,
}

\title{\textbf{LifeOS --- iOS Development Workflow}}
\author{}
\date{}

\begin{document}
\maketitle
\tableofcontents
\newpage

% -------------------------------------------------------
\section{Daily Development (JS/TS changes)}

\begin{enumerate}[leftmargin=*, label=\textbf{Step \arabic*.}]

  \item \textbf{Start the dev server.} Open a Terminal, navigate to the project, and run:
  \begin{lstlisting}
cd ~/Desktop/LifeApp
./dev.sh
  \end{lstlisting}
  The script auto-detects your Mac's LAN IP and prints the connection URL.
  Keep this terminal open the entire coding session.

  \item \textbf{Open the app on your iPhone} (wired to the Mac, screen unlocked).
  \begin{itemize}
    \item \textbf{Auto-connects:} nothing to do.
    \item \textbf{Shows ``No development servers found'':} tap \textbf{Enter URL manually}
      and type the \texttt{exp://} URL printed by the script
      (e.g.\ \texttt{exp://192.168.0.114:8081}).
      You only need to do this \emph{once} --- the app remembers the address.
  \end{itemize}

  \item \textbf{Make changes and save.}
  Fast Refresh pushes every \texttt{.ts}/\texttt{.tsx} save to the device
  instantly --- no manual reload required.

\end{enumerate}

\medskip
\noindent\textbf{VSCode shortcut:} \texttt{Cmd+Shift+P} $\to$
\textit{Tasks: Run Task} $\to$ \textbf{Start Dev Server}
(no terminal needed, output stays in the integrated panel).

% -------------------------------------------------------
\section{Native Rebuild (new packages / iOS changes)}

Required when you:
\begin{itemize}
  \item Install an npm package that contains native iOS code
  \item Change \texttt{app.json} or \texttt{app.config.ts}
  \item Edit anything inside the \texttt{ios/} folder
\end{itemize}

Stop the dev server (\texttt{Ctrl+C}), then run:
\begin{lstlisting}
./run.sh
\end{lstlisting}
The script reads the device UDID from \texttt{.env}, rebuilds the native app,
and installs it on the device. Start \texttt{./dev.sh} again afterwards.

% -------------------------------------------------------
\section{Release / App Store Build}

Produces a \texttt{Release} binary with no dev tools or Metro dependency:
\begin{lstlisting}
./release.sh
\end{lstlisting}

% -------------------------------------------------------
\section{Fast Refresh Limitations}

Fast Refresh works for the vast majority of edits. It will \emph{not} fully
apply in these cases --- shake the device and tap \textbf{Reload} instead:

\begin{itemize}
  \item Editing a file that exports a class component
  \item Changing the root navigator structure
  \item Modifying a custom hook that has side effects in its module scope
\end{itemize}

% -------------------------------------------------------
\section{Quick Reference}

\begin{center}
\begin{tabular}{>{\ttfamily}ll}
\toprule
\normalfont\textbf{Command} & \textbf{When to use} \\
\midrule
./dev.sh         & Every coding session (start here) \\
./run.sh         & After installing a native package or changing iOS config \\
./release.sh     & App Store / TestFlight submission \\
\bottomrule
\end{tabular}
\end{center}

\medskip
\noindent\textbf{Device UDID} is stored in \texttt{.env} (git-ignored).
To update it, edit that file.

\end{document}
